\section{Threats to Validity}
This study has limitations that must be considered when interpreting the results. In this section, we discuss these limitations by construct, internal and external validity.

\subsection{Construct Validity}
First, regarding construct validity, the experiments use fixed instruction configurations for each task performed by the LLMs, such as code generation, hallucination detection, and the interaction cycle. This choice is intentional to preserve comparability and avoid additional variation. In this way, the results characterize the model's behavior under this common instruction protocol and token set, instead of performing a prompt search that would optimize all models in general, or a specific model, but would add another body of variation that is outside the objective.

The proposed taxonomy of errors is a result of executing the generated code, encompassing both execution and test results, and can be observed through the four levels presented: \emph{Syntax Error}, \emph{Execution Error}, \emph{Functional Error}, and \emph{Correct}. However, these classifications are less detailed than those in related work, but they go beyond a binary evaluation of error and correctness, allowing broader analysis of the main errors found in the code. Furthermore, our taxonomy did not address macro-level issues already treated as non-functional requirements, such as security, performance, and maintainability, because these are not observable in the execution output.

Another taxonomy issue concerns the classification of syntax errors. Because the execution process is deterministic and fixed, malformed code, even if it deviates in ways that would cause a runtime or functional error, is marked only by the major error associated with the delusion of following the established grammar, since it is not compiled. Another relevant point about syntax execution is that, because it uses a single protocol for all languages, models with different response patterns produce text mixed with code or malformed code. We treated this as a delusion in our work, since the model deviated from the grammar requested in the prompt. This was especially evident in the Ornith 1.5 35B and Gemma 4 31B models: while the other SLMs had a syntax error rate of less than 1\%, these models reached 10.1\% and 3.2\%, respectively. In some cases, it is possible to extract the code, but to maintain the methodological rigor applied in all other experiments, we applied the error label.

Finally, the label \emph{Correct} is determined only from the available test sets of EvalPlus, not from a real proof that the model will have the expected behavior for all possible inputs, so it does not serve as a formal proof of semantic correctness.

\subsection{Internal Validity}

Another limitation concerns the models used in the comparisons. The local SLMs used for the tasks were mostly quantized. We adopted this approach based on prior empirical work showing no significant loss when quantizing models to 4 bits and that, with quantization, we could use larger models given hardware limitations. To maintain methodological rigor, we used all models in quantized form in the comparison, even those that would fit the computation without precision reduction. In one model used for comparison, we noted that the unquantized variant performed slightly better than the quantized version. This difference does not invalidate these models; however, model performance should not be interpreted as the total performance of the model without precision reduction. Furthermore, this result shows that larger quantized models may be more advantageous than smaller lossless models.

In our LLM experiments, we performed inference from a single generation rather than from multiple responses to the same input, since LLMs behave non-deterministically and responses may vary, although we applied procedures to stabilize their behavior. Therefore, future work could investigate generating multiple outputs for a given task and, through statistical selection, choose the most frequent outputs as the final result.

Another limitation concerns the models selected for specialization, since we evaluated only Gemma models that fit our hardware limitations. The improvements show that specialization can be effective in the evaluated scenario, but this pattern may not hold for other models or adaptation methods. Furthermore, limitations included the limited range of hyperparameters, the number of runs, and the optimization targets. Future work could expand the scope of hyperparameters by conducting more searches.

The distillation process can introduce noise into the supervision and should not be treated as entirely correct, as the model's explanations may ignore error cases in the code or present false or inaccurate information. However, the introduction of rationales also builds on prior work showing that rationales can improve the model's generalization, rather than simply making it predict a class.

Another limitation lies in the feedback given to the models, because, as seen in RQ4, Pyright is not a source of absolute truth like the language compiler, and therefore, it can provide feedback that may be detrimental to the model, as seen in the results of RQ5, where some models had lower performance in recalling the detection of the \emph{Runtime Error} class. However, even with a warning that the result would not necessarily be correct.

\subsection{External Validity}

The experiments are limited to Python and do not imply the same results for the same task in other languages. Furthermore, the benchmark used only contains isolated tasks at the function level and does not reflect more robust behavior in real-world programming scenarios, such as those involving repositories. This benchmark is used to conduct a study. Initially, to perform a direct comparison with works that have already mapped the main errors in the same tasks, to show the current state of the models and how the difference between larger and smaller models has changed throughout the analyses. Another motivation for using this dataset, rather than real repositories, is that models would need to perform additional skills, such as environment interactions and agent studies, during the execution cycle, which would fall outside the task's overall scope.

The nominal parameter count should be considered only as an indicator of difference, and not as the only absolute difference between the models. Although some models of the same version, with variations in size, may have similar internal factors, there are still other factors that differ between models, in addition to size, such as the training dataset, the internal architecture during all phases related to training, among other factors that cannot be controlled internally.